%Data institusi

%\input{Dozentdaten}

%\renewcommand{\fachbereich}{Fachbereich}

%\renewcommand{\dozent}{Prof. Dr. . }

%Data ujian

\renewcommand{\klausurgebiet}{ }

\renewcommand{\klausurtyp}{ }

\renewcommand{\klausurdatum}{ . 20}

\klausurvorspann {\fachbereich} {\klausurdatum} {\dozent} {\klausurgebiet} {\klausurtyp}

%Data untuk tabel nilai berikut

\renewcommand{\aeins}{  3  }

\renewcommand{\azwei}{  3   }

\renewcommand{\adrei}{  7  }

\renewcommand{\avier}{  4  }

\renewcommand{\afuenf}{  0  }

\renewcommand{\asechs}{  0  }

\renewcommand{\asieben}{  0  }

\renewcommand{\aacht}{  5  }

\renewcommand{\aneun}{  15  }

\renewcommand{\azehn}{  5  }

\renewcommand{\aelf}{  0  }

\renewcommand{\azwoelf}{  3  }

\renewcommand{\adreizehn}{  3   }

\renewcommand{\avierzehn}{  0  }

\renewcommand{\afuenfzehn}{  10  }

\renewcommand{\asechzehn}{  4  }

\renewcommand{\asiebzehn}{  62  }

\renewcommand{\aachtzehn}{    }

\renewcommand{\aneunzehn}{    }

\renewcommand{\azwanzig}{    }

\renewcommand{\aeinundzwanzig}{    }

\renewcommand{\azweiundzwanzig}{    }

\renewcommand{\adreiundzwanzig}{    }

\renewcommand{\avierundzwanzig}{    }

\renewcommand{\afuenfundzwanzig}{    }

\renewcommand{\asechsundzwanzig}{    }

\punktetabellesechzehn

\klausurnote

\newpage

\setcounter{section}{K}

\inputaufgabegibtloesung
{3}
{

Definisikan istilah-istilah berikut
\zusatzklammer {yang dicetak miring} {} {}.
\aufzaehlungsechs{Suatu kurva
\stichwort {berparametrisasi panjang busur} {}
\maabbdisp {\gamma} {[a,b]} { \R^n
} {.}

}{Untuk suatu titik

\mavergleichskette
{\vergleichskette
{  P
}
{ \in }{ M
}
{  }{}
{    }{}
{   }{}
}
{}{}{}
pada suatu manifold diferensiabel $M$, yang dimaksud dengan dua kurva diferensiabel yang
\stichwort {ekuivalen secara tangensial} {} ialah
\maabbdisp {\gamma_1, \gamma_2} { I } { M
} {}
\zusatzklammer {di sini

\mavergleichskettek
{\vergleichskettek
{  0
}
{ \in }{  I
}
{  }{
}
{    }{
}
{   }{
}
}
{}{}{}
merupakan interval riil terbuka dan

\mavergleichskettek
{\vergleichskettek
{   \gamma_1(0)
}
{ = }{ \gamma_2(0)
}
{ = }{ P
}
{    }{
}
{   }{
}
}
{}{}{}} {} {.}

}{Suatu
\stichwort {penampang kontinu} {}
dari suatu pemetaan kontinu
\maabbdisp {p} {Y} {X
} {}
antara ruang-ruang topologis
\mathkor {} {X} {dan} {Y} {.}

}{\stichwort {Bentuk volume kanonik} {} pada suatu manifold Riemann berorientasi $M$.

}{\stichwort {Turunan eksterior} {}
dari suatu
\definitionsverweis {bentuk diferensial}{}{}
berderajat $k$ yang
\definitionsverweis {terdiferensialkan secara kontinu}{}{}

\mathl{\omega \in
{ \mathcal E }^{ k } (  U   )}{} pada suatu
\definitionsverweis {himpunan terbuka}{}{}

\mathl{U \subseteq \R^n}{.}

}{\stichwort {Kelengkungan seksional} {} yang berkaitan dengan vektor-vektor bebas linear

\mavergleichskette
{\vergleichskette
{  v,w
}
{ \in }{T_PM
}
{  }{
}
{    }{
}
{   }{
}
}
{}{}{}
pada suatu
\definitionsverweis {manifold Riemann}{}{}
$M$ di suatu titik

\mavergleichskette
{\vergleichskette
{ P
}
{ \in }{ M
}
{  }{
}
{    }{
}
{   }{
}
}
{}{}{.}
}

}
{} {}

\inputaufgabegibtloesung
{3}
{

Rumuskan teorema-teorema berikut.
\aufzaehlungdrei{\stichwort {Teorema tentang orientasi pada suatu hipermuka diferensiabel} {}

\mavergleichskette
{\vergleichskette
{  Y
}
{ \subseteq }{  \R^n
}
{  }{
}
{    }{
}
{   }{
}
}
{}{}{.}}{Rumus braket Lie medan-medan vektor pada suatu himpunan terbuka

\mavergleichskette
{\vergleichskette
{  U
}
{ \subseteq }{  \R^n
}
{  }{
}
{    }{
}
{   }{
}
}
{}{}{.}}{Teorema tentang
\stichwort {retraksi ke batas} {}
pada manifold.}

}
{} {}

\inputaufgabegibtloesung
{7}
{

Buktikan
\stichwort {teorema tentang penyelesaian persamaan diferensial biasa pada suatu hipermuka diferensiabel} {}

\mavergleichskette
{\vergleichskette
{ Y
}
{ \subseteq }{ \R^n
}
{  }{
}
{    }{
}
{   }{
}
}
{}{}{.}

}
{} {}

\inputaufgabegibtloesung
{4}
{

Misalkan
\maabb {h} { I } { V
} {}
sebuah kurva yang terdiferensialkan dua kali dalam
\definitionsverweis {ruang vektor Euklides}{}{}
$V$. Buktikan bahwa jika

\mavergleichskette
{\vergleichskette
{h'(t)
}
{ \neq }{ 0
}
{  }{
}
{    }{
}
{   }{
}
}
{}{}{,}
maka berlaku persamaan

\mavergleichskettedisp
{\vergleichskette
{ \Vert { h'(t)} \Vert'
}
{ =} { { \frac{   \left\langle h'(t) , h^{\prime \prime} (t)  \right\rangle  }{  \left\langle h'(t) , h'(t)  \right\rangle  } } \Vert { h'(t)} \Vert
}
{ } {
}
{ } {
}
{ } {
}
}
{}{}{.}

}
{} {}

\inputaufgabe
{0}
{

}
{} {}

\inputaufgabe
{0}
{

}
{} {}

\inputaufgabe
{0}
{

}
{} {}

\inputaufgabegibtloesung
{5}
{

Buktikan pernyataan bahwa ekuivalensi tangensial lintasan-lintasan pada suatu manifold di suatu titik $P$ dapat diperiksa menggunakan sebarang peta.

}
{} {}

\inputaufgabegibtloesung
{15 (3+3+4+5)}
{

Misalkan

\mavergleichskette
{\vergleichskette
{T
}
{ = }{S^1 \times S^1
}
{  }{
}
{    }{
}
{   }{
}
}
{}{}{}
merupakan torus dan

\mavergleichskette
{\vergleichskette
{S^2
}
{ \subset }{\R^3
}
{  }{
}
{    }{
}
{   }{
}
}
{}{}{}
merupakan
\definitionsverweis {permukaan bola satuan}{}{.}

a) Tunjukkan bahwa
\maabbeledisp {\varphi} {S^1 \times S^1} { S^2
} { \begin{pmatrix}  \alpha \\ \beta   \end{pmatrix} } { { \frac{   1 }{  2 } } \begin{pmatrix}  \sqrt{ 2-2 \sin \alpha   } \cos \beta  \\ \cos \alpha + \sin \beta \cos \alpha \\  1  + \sin \alpha -  \sin \beta + \sin \alpha \sin \beta   \end{pmatrix}
} {}
mendefinisikan suatu pemetaan kontinu.

b) Tunjukkan bahwa $\varphi$ surjektif.

c) Deskripsikan serat-serat $\varphi$.

d) Jelaskan pemetaan $\varphi$ dengan bantuan sebuah sketsa.

}
{} {}

\inputaufgabegibtloesung
{5}
{

Hitunglah integral lintasan
\mathl{\int_\gamma \omega}{} untuk
\maabbeledisp {\gamma} { [-1,0] } {\R^3
} { t } {(-t^2,t^3-1,t+2)
} {,}
dengan bentuk diferensial berderajat $1$

\mavergleichskettedisp
{\vergleichskette
{ \omega
}
{ =} { x^3dx -yzdy +xz^2 dz
}
{ } {
}
{ } {
}
{ } {
}
}
{}{}{}
pada $\R^3$.

}
{} {}

\inputaufgabe
{0}
{

}
{} {}

\inputaufgabegibtloesung
{3 (2+1)}
{

Misalkan
\maabbeledisp {f} {\R} {\R
} { t } {f(t)
} {,}
diberikan oleh

\mavergleichskettedisp
{\vergleichskette
{ f(t)
}
{ =} { { \frac{   \sin^{ 3  } (t^4)  }{  1+t^2 } }
}
{ } {
}
{ } {
}
{ } {
}
}
{}{}{}

a) Hitung turunan eksterior dari $f$.

b) Hitung turunan eksterior dari $fdt$.

}
{} {}

\inputaufgabegibtloesung
{3}
{

Berikan suatu contoh pemetaan yang terdiferensialkan secara kontinu
\maabbdisp {\varphi} { H } { H
} {}
dari suatu setengah bidang Euklides ke dirinya sendiri, dengan sifat bahwa tepat satu titik batas $H$ dipetakan ke suatu titik batas dan semua titik lainnya dipetakan ke interior setengah bidang tersebut.

}
{} {}

\inputaufgabe
{0}
{

}
{} {}

\inputaufgabegibtloesung
{10}
{

Buktikan versi balok dari teorema Stokes.

}
{} {}

\inputaufgabegibtloesung
{4}
{

Untuk kurva
\maabbeledisp {\psi} {\R} { {\mathbb H}
} { t } { \left(  t   , \,  e^t                   \right)
} {,}
yang bernilai dalam
\definitionsverweis {setengah bidang Poincaré}{}{}
${\mathbb H}$, tentukan
\definitionsverweis {simbol-simbol Christoffel}{}{}
untuk
\definitionsverweis {koneksi tarikan balik}{}{}
\zusatzklammer {dari
\definitionsverweis {koneksi Levi-Civita}{}{}
pada ${\mathbb H}$} {} {}
di atas $\R$.

}
{} {}

 [[Kategorie:Latexseite]]
